
\subsubsection{Third Section}

The third section contains 6 pieces, 5 of which are editorial articles. We shall highlight the polemic against "Lutheraneren," Chairman Homme, Dr. Schmidt — "The Old Church Order," 1897-98 — so instructive in so many respects. Likewise: "Do You Want a Congregation?" in which Sverdrup crosses swords with T. N. S. Finally, "The Biblical Legitimacy of Our Congregations and the Relation of God’s Children to a Free-Church Congregation," 1898, in which Pastor Müller-Eggen’s new concept of the congregation, sanctioned as it was by The United Church, receives its condemnation in black and white.

Here, then, that man is the critic, concerning whom the review column of "Lutheraneren" recently said: "As clear as Sverdrup’s view is when it concerns the invisible Church, so hazy does it become when there is question of the visible congregation and the relation between the two."

We must ask: "Whose view was hazy? Was it Sverdrup’s or Müller-Eggen’s, or theirs who sanctioned Eggen’s concept of the congregation, including the clever comparison of visible congregations to barrels of salted meat?"

The fourth piece in this section, "Can One Minister to a Believer?" is an assessment of Pastor Lundeberg’s congregational program.

\subsubsection{Fourth Section}

The whole fourth section is devoted to lay activity. In one essay the memory of Hauge’s calling as a preacher of repentance is celebrated. Another treats T. A. S.’s theory of evolution as applied to lay activity. A third discusses lay activity as it was treated in 1906 by "the great Joint Committee which is negotiating concerning the union of the church bodies."

\subsubsection{Fifth Section}

The fifth section, in three articles, treats infant baptism and confirmation.

\subsubsection{Sixth Section}

The sixth section contains lectures on congregational life, the condition for the formation of a congregation, the significance of mission for the congregation, and the unity of believers.

These lectures are, if one may so speak, somewhat abstract by comparison with the fuller contributions of earlier days. But it is always the same goal, the same principles, that are at issue.

We break off. From these scattered remarks one will perhaps gain some little idea of the abundance of material which the book presents. The work itself offers the best survey through the introductions and notes with which its editor, Professor Helland, has furnished the individual sections and articles.

In his excellent essay "On the Conditions for a Norwegian Theology" (1908), Olaf Moe, docent at the University of Christiania, writes: "Thus the chapter on the Church also becomes one of the weakest in our practical dogmatics. One understands easily that the Church is invisible, but one understands only poorly that it is also visible; one perhaps believes in the forgiveness of sins, but does not know that it is ‘in this Christian congregation we receive a daily full forgiveness of sins.’" When a man such as Docent Moe can write thus — and who can deny that he has hit the mark? — it may be hoped that Sverdrup’s book on the congregation will have its importance for the Norwegian churchfolk on both sides of the ocean. — The third volume, which will soon be published, will be discussed later.


\subsection {Advertisements included in the book} 

The first three volumes of Professor Georg Sverdrup’s Selected Writings have been published

Vol. 1:
Sketches and Essays
400 pages

Vol. 2:
On the Congregation
395 pages

Vol. 3:
On Augsburg and the Free Church
394 pages

The whole work will comprise 5 to 6 volumes, and the remaining ones will be published at intervals of 6 months and will comprise the same number of pages each as those already published.

Each volume costs:
In half Morocco .......... $2.50
In Silk Cloth ............ $1.80

Sold either by subscription to the complete work, or in separate volumes.


Agents wanted.
Liberal commission offered.
Write for terms.

The Free Church Book Concern,
322 Cedar Ave.,
Minneapolis, Minn.
